\subsection{恒电流产生的磁场 (稳磁场)}
稳磁场问题的基本方程为
\begin{equation}
    \nabla\cdot\bm{B}=0.
\end{equation}

容易得到 $\bm{B}$ 可以表示为
\begin{equation}
    \bm{B}=\nabla\times\bm{A}.
\end{equation}
其中 $\bm{A}$ 是\textbf{磁矢位}.

而
\begin{equation}
    \nabla\times\bm{H}=\frac{1}{\mu}[\nabla(\nabla\cdot\bm{A})-\nabla^2\bm{A}].
\end{equation}
令磁场满足\textbf{库仑规范}: $\nabla\cdot\bm{A}=0$, 则
\begin{equation}
    \nabla\times\bm{H}=-\frac{1}{\mu}\nabla^2\bm{A}=\bm{J}.
\end{equation}

因此
\begin{equation}
    \nabla^2\bm{A}=-\mu\bm{J}.
\end{equation}

这是三个方向上的泊松方程. 给定边界条件之后, 即可得到唯一解. 通解表为
\begin{equation}
    \bm{A}(\bm{r})=-\frac{\mu}{4\pi}\int\frac{\bm{J}(\bm{r'})}{|\bm{r}-\bm{r'}|}\mathrm{d}v'.
\end{equation}
边界条件表为
\begin{equation}
    \begin{cases}
        \bm{A_1}=\bm{A_2}, \\
        \bm{e_n}\times\left(\frac{1}{\mu_1}\nabla\times\bm{A_1}-\frac{1}{\mu_2}\nabla\times\bm{A_2}\right)=\bm{J_S}.
    \end{cases}
\end{equation}

考虑柱坐标:
\begin{equation}
    \bm{A}=-\frac{\mu}{4\pi}\int_{0}^{2\pi}\frac{\bm{J}(\bm{r'})}{|\bm{r}-\bm{r'}|}a\mathrm{d}\varphi.
\end{equation}
其中
\begin{gather*}
    \bm{r}=\bm{e_r}r=\bm{e_x}r\sin\theta+\bm{e_z}r\cos\theta, \\
    \bm{r'}=\bm{e_x}a\cos\varphi'+\bm{e_y}a\sin\varphi, \\
    \bm{J}=i\mathrm{d}\bm{l}, \\
    \mathrm{d}\bm{l}=\bm{e_{\varphi'}}a\mathrm{d}\varphi'=-\bm{e_x}\sin\varphi'+\bm{e_y}\cos\varphi', \\
    |\bm{r}-\bm{r'}|=\sqrt{r^2+a^2-2ar\sin\theta\cos\varphi'}, \\
    \frac{1}{|\bm{r}-\bm{r'}|}\approx\frac{1}{r}\left(1+\frac{d}{r}\sin\theta\cos\varphi'\right).
\end{gather*}
可以得到
\begin{equation}
    \bm{A}=\bm{e_y}\frac{\mu\pi ia^2}{4\pi r^2}\sin\theta. \qquad (\bm{e_y}=\bm{e_\varphi})
\end{equation}
